\subsection{অক্টাল থেকে বাইনারি}
প্রতি অক্টাল অঙ্কে ৩ বিট, প্রতি হেক্সাডেসিমাল অঙ্কে ৪ বিট। প্রতিটি অঙ্কের বিটগুলো একই ক্রমে একত্র করো; ভগ্নাংশের শুরুর শূন্য বাদ দিও না।
\par\noindent\begin{minipage}{\linewidth}
\convExample{১}{পূর্ণসংখ্যা}
$ (265)_{8} $-কে ২ ভিত্তিতে রূপান্তর করো।\par
\par\noindent\begin{minipage}{\linewidth}\centering
\phantomsection\label{VC036}
\begin{chThreeDiagram}\begin{tikzpicture}[x=1mm,y=1mm]
\node[vc digit] (s0) at (6.0,0) {$ 2 $};
\draw[vc arrow] (s0.south)--(6.0,-9);
\node[vc digit] (s1) at (24.0,0) {$ 6 $};
\draw[vc arrow] (s1.south)--(24.0,-9);
\node[vc digit] (s2) at (42.0,0) {$ 5 $};
\draw[vc arrow] (s2.south)--(42.0,-9);
\node[vc bit] (raw0) at (0,-12) {$ 0 $};
\node[vc bit] (raw1) at (6,-12) {$ 1 $};
\node[vc bit] (raw2) at (12,-12) {$ 0 $};
\node[vc bit] (raw3) at (18,-12) {$ 1 $};
\node[vc bit] (raw4) at (24,-12) {$ 1 $};
\node[vc bit] (raw5) at (30,-12) {$ 0 $};
\node[vc bit] (raw6) at (36,-12) {$ 1 $};
\node[vc bit] (raw7) at (42,-12) {$ 0 $};
\node[vc bit] (raw8) at (48,-12) {$ 1 $};
\draw[vc brace] (-3.0,-16)--(15.0,-16);
\draw[vc brace] (15.0,-16)--(33.0,-16);
\draw[vc brace] (33.0,-16)--(51.0,-16);
\end{tikzpicture}\end{chThreeDiagram}
\end{minipage}\par
\[\therefore\ (265)_{8}=(10110101)_{2}.\]
পূর্ণাংশের শুরুর বা ভগ্নাংশের শেষের অপ্রয়োজনীয় ০ বাদ দিয়ে ফল লেখা হয়েছে; ভগ্নাংশের শুরুর ০ বাদ যায় না।\par
\end{minipage}\par
\par\noindent\begin{minipage}{\linewidth}
\convExample{২}{শুধু ভগ্নাংশ}
$ (0.47)_{8} $-কে ২ ভিত্তিতে রূপান্তর করো।\par
\par\noindent\begin{minipage}{\linewidth}\centering
\phantomsection\label{VC037}
\begin{chThreeDiagram}\begin{tikzpicture}[x=1mm,y=1mm]
\node[vc digit] (s0) at (6.0,0) {$ 0 $};
\draw[vc arrow] (s0.south)--(6.0,-9);
\node[vc digit] (s1) at (36.0,0) {$ 4 $};
\draw[vc arrow] (s1.south)--(36.0,-9);
\node[vc digit] (s2) at (54.0,0) {$ 7 $};
\draw[vc arrow] (s2.south)--(54.0,-9);
\node[vc digit] at (21.0,0) {$.$};
\node[vc bit] (raw0) at (0,-12) {$ 0 $};
\node[vc bit] (raw1) at (6,-12) {$ 0 $};
\node[vc bit] (raw2) at (12,-12) {$ 0 $};
\node[vc bit] (raw3) at (30.0,-12) {$ 1 $};
\node[vc bit] (raw4) at (36.0,-12) {$ 0 $};
\node[vc bit] (raw5) at (42.0,-12) {$ 0 $};
\node[vc bit] (raw6) at (48.0,-12) {$ 1 $};
\node[vc bit] (raw7) at (54.0,-12) {$ 1 $};
\node[vc bit] (raw8) at (60.0,-12) {$ 1 $};
\node[vc digit] at (21.0,-12) {$.$};
\draw[vc brace] (-3.0,-16)--(15.0,-16);
\draw[vc brace] (27.0,-16)--(45.0,-16);
\draw[vc brace] (45.0,-16)--(63.0,-16);
\end{tikzpicture}\end{chThreeDiagram}
\end{minipage}\par
\[\therefore\ (0.47)_{8}=(0.100111)_{2}.\]
পূর্ণাংশের শুরুর বা ভগ্নাংশের শেষের অপ্রয়োজনীয় ০ বাদ দিয়ে ফল লেখা হয়েছে; ভগ্নাংশের শুরুর ০ বাদ যায় না।\par
\end{minipage}\par
\par\noindent\begin{minipage}{\linewidth}
\convExample{৩}{মিশ্র সংখ্যা}
$ (143.52)_{8} $-কে ২ ভিত্তিতে রূপান্তর করো।\par
\par\noindent\begin{minipage}{\linewidth}\centering
\phantomsection\label{VC038}
\begin{chThreeDiagram}\begin{tikzpicture}[x=1mm,y=1mm]
\node[vc digit] (s0) at (6.0,0) {$ 1 $};
\draw[vc arrow] (s0.south)--(6.0,-9);
\node[vc digit] (s1) at (24.0,0) {$ 4 $};
\draw[vc arrow] (s1.south)--(24.0,-9);
\node[vc digit] (s2) at (42.0,0) {$ 3 $};
\draw[vc arrow] (s2.south)--(42.0,-9);
\node[vc digit] (s3) at (72.0,0) {$ 5 $};
\draw[vc arrow] (s3.south)--(72.0,-9);
\node[vc digit] (s4) at (90.0,0) {$ 2 $};
\draw[vc arrow] (s4.south)--(90.0,-9);
\node[vc digit] at (57.0,0) {$.$};
\node[vc bit] (raw0) at (0,-12) {$ 0 $};
\node[vc bit] (raw1) at (6,-12) {$ 0 $};
\node[vc bit] (raw2) at (12,-12) {$ 1 $};
\node[vc bit] (raw3) at (18,-12) {$ 1 $};
\node[vc bit] (raw4) at (24,-12) {$ 0 $};
\node[vc bit] (raw5) at (30,-12) {$ 0 $};
\node[vc bit] (raw6) at (36,-12) {$ 0 $};
\node[vc bit] (raw7) at (42,-12) {$ 1 $};
\node[vc bit] (raw8) at (48,-12) {$ 1 $};
\node[vc bit] (raw9) at (66.0,-12) {$ 1 $};
\node[vc bit] (raw10) at (72.0,-12) {$ 0 $};
\node[vc bit] (raw11) at (78.0,-12) {$ 1 $};
\node[vc bit] (raw12) at (84.0,-12) {$ 0 $};
\node[vc bit] (raw13) at (90.0,-12) {$ 1 $};
\node[vc bit] (raw14) at (96.0,-12) {$ 0 $};
\node[vc digit] at (57.0,-12) {$.$};
\draw[vc brace] (-3.0,-16)--(15.0,-16);
\draw[vc brace] (15.0,-16)--(33.0,-16);
\draw[vc brace] (33.0,-16)--(51.0,-16);
\draw[vc brace] (63.0,-16)--(81.0,-16);
\draw[vc brace] (81.0,-16)--(99.0,-16);
\end{tikzpicture}\end{chThreeDiagram}
\end{minipage}\par
\[\therefore\ (143.52)_{8}=(1100011.10101)_{2}.\]
পূর্ণাংশের শুরুর বা ভগ্নাংশের শেষের অপ্রয়োজনীয় ০ বাদ দিয়ে ফল লেখা হয়েছে; ভগ্নাংশের শুরুর ০ বাদ যায় না।\par
\end{minipage}\par
